\section{Method}\label{sec:method}

\subsection{Prompt-conditioned calibration}

Let $z$ be the raw classification logit of a detection and $\phi$ the phrase that was used to
request its category. A scalar calibrator fits one map $z\mapsto\sigma(az+b)$ for all
detections. We instead let the two parameters be functions of the phrase:

\begin{equation}\label{eq:phc}
p = \sigma\!\left(z\,e^{-\ell(\phi)} + m(\phi)\right),
\qquad
\ell(\phi) = \log T(\phi),\quad m(\phi) = b(\phi),
\end{equation}

where $\ell$ and $m$ are produced by a linear head on a descriptor $x(\phi)$ of the phrase.
Since $T(\phi)=e^{\ell(\phi)}>0$ for every phrase, the map is strictly increasing in $z$ for
fixed $\phi$.

\begin{proposition}[The map cannot reorder a run]\label{prop:mono}
For any two detections that were requested with the same phrase, $z_i>z_j$ implies
$p_i>p_j$ under \eqref{eq:phc}. Consequently every ranking metric computed within a
single-wording run, including average precision and the risk--coverage curve, is unchanged by
the calibrator.
\end{proposition}

\begin{proof}
$T(\phi)>0$, so $z\mapsto z/T(\phi)+m(\phi)$ is strictly increasing, and $\sigma$ is strictly
increasing.
\end{proof}

This property is what makes the method safe to use. The map can change the score scale, which
is what thresholds and calibration error are about, and it cannot change the order of the boxes
that a user sees in one run, which is what AP is about. A map that may depend on $z$ in a
non-monotone way, such as isotonic regression, does not have this property, and
Table~\ref{tab:main} reports the AP that it costs.

\paragraph{The phrase descriptor.} $x(\phi)$ has 33 entries. The first 24 are a principal
component projection of the pooled last-layer text embedding that the detector's own text
encoder produced for $\phi$, averaged over tokens. This requires no extra forward pass, because
the encoder has already run. The remaining entries are the token count, the character length,
the embedding norm, the cosine similarity to the bare name of the same category, and a one-hot
encoding of the five wordings. The head is linear, so the whole calibrator has 68 parameters. We
use a linear head because the calibrator has to be fitted on the phrases that one is willing to
enumerate, and a linear map on interpretable descriptors degrades gracefully when the deployed
phrase lies outside that set.

\paragraph{Fitting.} The parameters minimise a weighted Brier loss on cached logits,
\begin{equation}
\mathcal{L} = \frac{1}{n}\sum_i w_i\,\big(p_i - y_i\big)^2 ,
\qquad
w_i = \begin{cases} w_{\rm band} & \text{if the uncalibrated score is in the band}\\ 1 & \text{otherwise,}\end{cases}
\end{equation}
with $w_{\rm band}=10$ in the default configuration. The band is computed once from the
uncalibrated score and then frozen. If it were recomputed from the output of the model, the
cheapest way to lower a band-weighted loss would be to push boxes out of the band rather than to
calibrate them.

\subsection{From a calibrated score to a threshold}

A calibrated probability is not yet a decision. The user of an annotation tool sets a threshold
$\lambda$ and keeps every box with $p\geq\lambda$. What the user cares about is the false
discovery rate among the kept boxes, $\mathrm{FDR}=\mathbb{E}[V/R]$, where $R$ is the number
kept and $V$ is the number of these that are wrong.

We choose $\lambda$ on a labelled calibration half by conformalised selection
\citep{jin2023selection}. The calibration detections are ordered by decreasing score. Let $R(k)$
and $V(k)$ be the number kept and the number wrong among the top $k$, and set

\begin{equation}\label{eq:lambda}
\hat\lambda = \min\Big\{\lambda \;:\; \frac{1+V(\lambda)}{\max(R(\lambda),1)} \leq \alpha\Big\},
\end{equation}

the smallest threshold whose conformal estimate of the false discovery proportion is below the
target $\alpha$. Under exchangeability of calibration and test detections this controls
$\mathbb{E}[\mathrm{FDR}]$ at $\alpha$ up to a term of order $1/R$
\citep{jin2023selection,angelopoulos2023conformal}.

The exchangeability assumption is the part that matters here. It holds comfortably when the
calibration and the test detections come from the same phrasing, and it fails when they do not.
A threshold calibrated on ``dog'' and applied to ``dogs'' is not covered by the theorem, and
Section~\ref{sec:exp-fdr} measures how far it misses.

One might expect that conditioning the score on the phrase restores the assumption, by making
detections requested with different wordings comparable. It does not, and
Proposition~\ref{prop:mono} shows why. The false discovery rate at a shared threshold is a
function of the \emph{order} of the kept boxes, and a monotone map leaves that order untouched.
If two wordings disagree about which boxes are the good ones, calibration cannot reconcile
them. What does work is to require the constraint within each wording:

\begin{equation}\label{eq:lambda-worst}
\hat\lambda_{\max}(\alpha) \;=\; \max_{v\in\{1..5\}} \hat\lambda_v(\alpha),
\end{equation}

where $\hat\lambda_v$ is \eqref{eq:lambda} computed on the calibration detections of wording $v$
alone. Since $\mathrm{FDR}$ falls as the threshold rises, \eqref{eq:lambda-worst} is the
strictest admissible choice, and it satisfies the constraint in every wording at the same time.
Because it is a maximum over strata, it needs no multiplicity correction. It is also the only one
of the three threshold protocols compared in Table~\ref{tab:fdr} that a tool can deploy as a
single number and still satisfy the constraint in every wording. Section~\ref{sec:exp-fdr}
reports what it costs.

\begin{table}[t]
\caption{The three protocols compared for every calibrator. A run is one phrasing on all
images. ``Mixed'' fits on all five wordings; ``held-out'' fits on three and is applied to the
two it never saw; ``transfer'' fits on COCO and is applied to LVIS.}\label{tab:protocols}
\begin{tabular*}{\textwidth}{@{\extracolsep{\fill}}lll@{}}
\toprule
Protocol & Fitted on & Evaluated on \\
\midrule
mixed & five wordings, half the images & the other half, all five \\
held-out wording & bare, article, photo & image, plural \\
transfer & COCO, five wordings & LVIS, five wordings, no refit \\
\bottomrule
\end{tabular*}
\end{table}
